%% ch16 --- Chaos inside the computer
%%
%% Written September 2026. Every chaotic picture in the book was drawn by a
%% machine that rounds at every step; this chapter asks what that does, and
%% measures it. Floating point from scratch; the doubling map dying on a
%% computer; one rule written two ways parting company within about fifty
%% steps, against a truth carried to a hundred digits; why exact and
%% extended arithmetic cannot rescue a long chaotic run; the rescue that does
%% work (shadowing, stated in a rigourbox and illustrated by finding a shadow
%% and by statistics that agree); and finite machines, loops and
%% random-number generators. Every number comes from code/measure/ch16.py,
%% which imports the same functions the listings print.
%%
%% WHAT THIS CHAPTER DOES NOT SAY. The Lyapunov exponent and the horizon
%% formula are Chapter 7's; they are used here, not derived. Lorenz's 1961
%% rounded restart is Chapter 6's story and is not retold. The Jacobian
%% products of machine learning, named in the manifest brief, are left out:
%% the chapter's word budget went to the measurements instead.
%%
%% Twin: chapters/pl/ch16-computers.tex. Section for section, macro for
%% macro, number for number; tools/parity.py compares them.

\chapter{Chaos inside the computer}
\label{ch:computers}
\index{floating point}\index{rounding}

\noindent
Every picture of chaos in this book was drawn by a computer, and so was
every number: the two runs peeling apart, the butterfly-shaped attractor,
the Lyapunov exponent to three decimal places. Here is an awkward fact about
the machine that drew them. It cannot store the number $\num{0.1}$. It
stores a number very close to it, and every time it adds or multiplies two
numbers it rounds the answer to one it can store.

You have spent several chapters learning that a chaotic system turns the
smallest error into a large one. Put the two facts together. If the computer
makes a tiny error at every step, and chaos magnifies every error, how can
any of those pictures be right? This chapter finds out, by measuring, and
the answer is better than you might fear and worse than you might hope.

\begin{outcomes}
\outcome{say what a computer stores when you type a number, and why
  $\num{0.1} + \num{0.2}$ does not come out as $\num{0.3}$}
\outcome{predict the step at which the doubling map dies on a computer, from
  the number it starts with}
\outcome{estimate how many steps a chaotic computation stays right from the
  number of binary digits it keeps}
\outcome{say which results of a chaotic simulation you can trust, and why}
\outcome{describe a random-number generator as a deterministic rule, and say
  why it must repeat}
\end{outcomes}

\section{What a computer keeps}
\label{sec:ch16-keeps}
\index{binary numbers}\index{double precision}

\noindent
Write one third as a decimal and you never finish: $\num{0.333}\ldots$, the
threes running on for ever. Any fixed number of decimal places gives a
number close to one third and not equal to it. Decimals can write down
exactly only the fractions whose bottom, in lowest terms, is built from twos
and fives, the factors of ten.

A computer counts in twos, not tens. In \emph{binary} the places after the
point are worth a half, a quarter, an eighth, a sixteenth and so on, so
$\tfrac{3}{4}$ is written $\num{0.11}$: a half plus a quarter.

\begin{think}
Which of these can binary write exactly, with a fixed number of places:
$\tfrac{1}{2}$, $\tfrac{3}{8}$, $\tfrac{1}{10}$?
\end{think}
\reveal{$\tfrac{1}{2}$ and $\tfrac{3}{8}$, but not $\tfrac{1}{10}$. Binary
writes exactly the fractions whose bottom is a power of two, and ten has a
factor of five.}
In binary, one tenth is $\num{0.000110011}\ldots$, with the block $0011$
repeating for ever, just as the threes repeat in one third. So a computer
has to cut it off somewhere. Python, like almost every program you use,
keeps a number as a \emph{double}: \val{ch16.bits.double} significant binary
digits, which is about \val{ch16.decimal.double} significant decimal digits,
together with a power of two that says where the point goes. For one tenth
it keeps the nearest double there is, which is too large by about
\val{ch16.tenth.err}. That is the first rounding. Every sum and every
product is rounded again: the exact answer is worked out and replaced by the
nearest double.

\begin{think}
Both $\num{0.1}$ and $\num{0.2}$ are stored slightly wrong, and their sum is
rounded once more. Does the computer agree that $\num{0.1} + \num{0.2}$ is
$\num{0.3}$?
\end{think}
\reveal{No. The sum prints as a number a whisker above $\num{0.3}$, and the
comparison answers \code{False}.}

\pyfile{code/ch16/floats.py}{What the computer stores, and what it does with it}{lst:ch16-floats}

\transcript{ch16-floats}

\noindent
The third line is the exact value the computer keeps when you type
$\num{0.1}$. The fourth writes the same number as a fraction, and the fifth
shows that the bottom of that fraction is $2^{\val{ch16.tenth.power}}$.
Every double is like that: a whole number divided by a power of two. The
sixth line adds up three numbers in two different orders and gets two
different answers, because each order rounds at different moments \dash{}
the grouping that school told you cannot matter does matter once every step
rounds. And the last line shows the room running out: past $2^{53}$ a double
cannot even hold every whole number.

\begin{note}
The rules for this rounding are the same on almost every computer you will
meet. They were fixed in 1985 in a standard called IEEE~754, which requires
$+$, $-$, $\times$, $\div$ and the square root to give the exact answer
rounded to the nearest double. That promise is why this book can print
numbers from chaotic calculations and expect your computer to agree with them
to the last digit: its own measurements use only those operations, and are
rerun and compared, digit for digit, every time the book is built.
\end{note}

\section{A map the computer cannot run}
\label{sec:ch16-doubling}
\index{doubling map}

\noindent
The simplest chaotic rule there is fits in a line: double the number, and if
the answer is $1$ or more, take $1$ away. In symbols, $x \to 2x \bmod 1$,
read \enquote{$2x$ with its whole part thrown away}. This is the
\emph{doubling map}. It stretches the gap between any two nearby starts by a
factor of two at every step, so its errors double exactly as the logistic
map's do on average at $r = 4$.

In binary you can watch it work. Doubling moves every binary digit one place
to the left, and taking away $1$ throws away the digit that has crossed the
point: $\num{0.1011}$ becomes $\num{0.011}$, then $\num{0.11}$, then
$\num{0.1}$. The map reads out its start one binary digit at a time.

\begin{think}
Work the doubling map by hand from $\tfrac{1}{10}$, in fractions, for five
steps. What does its orbit do for ever after?
\end{think}
\reveal{$\tfrac{1}{5}$, $\tfrac{2}{5}$, $\tfrac{4}{5}$, $\tfrac{3}{5}$,
$\tfrac{1}{5}$: from the first step on it goes round the same four values for
ever, and never reaches $0$.}
Now think about the computer's version. Its $\num{0.1}$ is a whole number
divided by $2^{\val{ch16.tenth.power}}$, so it has a last binary digit, and
every step pushes one digit out past the point.

\begin{think}
What will the computer's orbit of $\num{0.1}$ be after sixty steps?
\end{think}
\reveal{Exactly $0$. When the last binary digit leaves, at step
\val{ch16.collapse.tenth}, nothing is left, and $0$ doubled is $0$ for ever.}

\pyregion{code/ch16/doubling.py}{race}{The doubling map from one tenth, on the computer and in exact fractions}{lst:ch16-doubling}

\transcript{ch16-doubling}

\noindent
For forty steps the computer and the truth agree to six decimal places. Then
the computer drifts, and at step \val{ch16.collapse.tenth} it dies, while the
true orbit goes on round its loop of four. A computed orbit can be wrong not
merely in its digits but in its whole character: a loop has become a fall to
nothing.

Here is the strange part. The computer's arithmetic made no error at all:
doubling a double and taking away $1$ never needs rounding, and the chapter's
measurement script checks, step by step, that the computer's orbit is exactly
the true orbit of the number it stored. What is wrong is the start. Every
number a computer can store has finitely many binary digits, and every such
start is swallowed by $0$. Almost every true start between $0$ and $1$ has
infinitely many digits and is never swallowed. The machine can hold only the
rare starts that die. From each of the \val{ch16.collapse.starts} starts
$\tfrac{1}{997}$, $\tfrac{2}{997}$, and so on, the computer's orbit reached
$0$ in between \val{ch16.collapse.min} and \val{ch16.collapse.max} steps, most
often in \val{ch16.collapse.mode}.

\plotfig{ch16-collapse}{The doubling map from one tenth. The truth (open
circles) goes round a loop of four values for ever. The computer (dots)
agrees with it for about forty steps, drifts, and at step
\val{ch16.collapse.tenth} falls to exactly $0$, where it stays.}{the doubling
map dying on a computer}

\begin{lab}{l16_01_collapse}{When does the doubling map die?}
Write the doubling map, a function that runs it from a start and counts the
steps until the orbit is exactly $0$, and a function that predicts that count
without running anything. For the prediction, Python's
\code{x.as\_integer\_ratio()} gives the stored number as a whole number over a
power of two. Your two answers must agree for every start the test tries.
\end{lab}

\section{One rule, two programs}
\label{sec:ch16-formulas}
\index{logistic map!on a computer}\index{reproducibility}

\noindent
The doubling map is a special case: it never rounds, and it dies because its
starts are special. The logistic map of Chapter~\ref{ch:logistic} is more
typical, because nearly every step rounds. Here it is at $r = 4$, written two
ways.

\pyregion{code/ch16/two_formulas.py}{rules}{The logistic map at r = 4, written two ways}{lst:ch16-rules}

\noindent
Multiply out the bracket in the first and you get the second:
$4x(1 - x) = 4x - 4x^{2}$. On paper they are one rule.

\begin{think}
Start both from $\num{0.1}$ on the same computer and run them for seventy
steps. Will the two programs print the same numbers?
\end{think}
\reveal{No. They differ in the last binary digit from step
\val{ch16.formulas.first}, and by step \val{ch16.formulas.apart} they are
more than $\num{0.1}$ apart: two unrelated sequences.}
The two formulas round at different moments. The first rounds $1 - x$ and
then a product; the second rounds $4x^{2}$ and then a difference. Each error
is about one part in $10^{16}$, and the map doubles it at every step. In the
listing's output the right-hand column is the truth: the orbit of exactly
one tenth, carried with a hundred significant digits, which is far more than
seventy steps can use up.

\transcript{ch16-two-formulas}

\noindent
Which formula is right? For forty steps, both, to six decimal places. By step
\val{ch16.wrong.a} the first has left the truth, and by step
\val{ch16.wrong.b} so has the second. Neither formula is wrong; each is right
for about fifty steps, and after that it is following a different orbit.

\begin{trapbox}
\textbf{\enquote{A deterministic program always gives the same answer.}} Run
the same program on the same machine and it does, to the last digit: nothing
in it is random. But change the computer, the compiler, the version of a
maths library or the order in which a long sum is added up, and you change
which roundings happen \dash{} which is exactly what the two formulas above
did. Compilers are allowed, when asked for speed, to rearrange arithmetic or
to fuse a multiplication and an addition into one step with one rounding. In
most programs the difference stays in the last digit. In a chaotic one it
reaches the first digit in about fifty steps.
\end{trapbox}

\noindent
Fifty is not an accident, and you can predict it.

\begin{think}
At $r = 4$ an error roughly doubles at every step, as Chapter~\ref{ch:butterfly}
measured. A double's rounding error is about $10^{-16}$, and ten doublings
multiply an error by about a thousand. Roughly how many steps until an error
of $10^{-16}$ has grown to $\num{0.1}$?
\end{think}
\reveal{About fifty. From $10^{-16}$ to $10^{-1}$ is fifteen factors of ten,
and every three factors of ten take about ten doublings.
Chapter~\ref{ch:lyapunov}'s horizon, with the rounding error as the error in
the start, gives \val{ch16.predict.double}.}
Said in one sentence: at $r = 4$ the map uses up one binary digit of accuracy
per step, so a calculation that keeps \val{ch16.bits.double} binary digits
stays right for about fifty steps. \emph{Single precision}, the format
graphics cards and much machine-learning software use because it is twice as
compact, keeps \val{ch16.bits.single} binary digits, and the same arithmetic
predicts about \val{ch16.predict.single} steps. Measured, it lasts
\val{ch16.wrong.single}.

\plotfig{ch16-parting}{Distance from the true orbit of one tenth, on a
logarithmic scale, for the logistic map at $r = 4$. Each run's error doubles
every step, along the dotted line, until it is as big as the whole interval:
after about twenty steps in single precision and about fifty in
double.}{two precisions parting company}

\begin{lab}{l16_02_single}{Single against double}
Python's numbers are doubles, but you can round one to the nearest single by
packing it into four bytes and reading it back; the starter shows how. Write
the logistic map at $r = 4$ with every operation rounded to single precision,
run it beside the double version from the same start, and find the first step
at which the two differ by more than $\num{0.1}$. Compare it with the
prediction above.
\end{lab}

\begin{worldbox}
Machine learning meets this every day. A neural network is trained by adding
up millions of small numbers over and over on a graphics card, whose
thousands of parallel threads can finish in a different order each time, so a
sum is grouped differently and rounds differently from run to run. Training
is sensitive enough that two runs of the same program from the same random
seed can end with networks whose weights differ throughout \dash{} and which
usually score about the same on a test. The documentation of PyTorch, a
widely used library, says plainly that fully reproducible results are not
guaranteed across platforms and releases, and offers a switch that forces a
fixed order of operations, at a cost in speed.
\end{worldbox}

\section{Why more digits cannot save you}
\label{sec:ch16-exact}
\index{exact arithmetic}\index{fractions!exact}

\noindent
The obvious cure is to stop rounding. Python's \pkg{fractions} module keeps
every number as an exact fraction, top and bottom, however large they grow.
One tenth stays one tenth, and $4x(1 - x)$ from it gives $\tfrac{9}{25}$,
then $\tfrac{576}{625}$, with no error at all.

\begin{think}
From $\tfrac{9}{25}$ to $\tfrac{576}{625}$ the bottom of the fraction went
from $25$ to $625 = 25 \times 25$. If it is squared at every step, what
happens to the number of digits it has?
\end{think}
\reveal{It doubles, roughly, at every step: squaring a number about doubles
its number of digits.}

\pyregion{code/ch16/more_digits.py}{exact}{Exact fractions: the bottom is squared at every step}{lst:ch16-exact}

\transcript{ch16-more-digits}

\noindent
After ten steps the bottom of the fraction has \val{ch16.exact.ten} digits.
After fifty it would have about \val{ch16.exact.fifty}, and merely storing it
would take some \val{ch16.exact.fifty.tb} terabytes. Exact arithmetic gets
every step right and cannot finish the calculation.

The last four lines of the output try the practical cure: keep a fixed but
larger number of digits (Python's \pkg{decimal} module lets you choose how
many) and compare each run with a thousand-digit one. Sixteen digits last
\val{ch16.right.sixteen} steps, like a double. Twice the digits last twice
as long, and every decimal digit buys about \val{ch16.right.per.digit} more
steps. That is the one-binary-digit-per-step rule again, because a decimal
digit is worth about \val{ch16.right.per.digit} binary ones.

\begin{trapbox}
\textbf{\enquote{More decimal places make a simulation exact.}} They make it
right for longer, by a fixed number of steps for each digit, and never for
ever. To follow a chaotic orbit for a thousand steps you must carry about a
thousand binary digits, some three hundred decimal ones; for a million steps,
a million. And the start would have to be known that precisely too, which no
measurement will ever give you. This is Chapter~\ref{ch:lyapunov}'s cruel
arithmetic, met from the computer's side: precision buys a horizon that grows
by a constant, not in proportion.
\end{trapbox}

\section{Wrong path, right picture}
\label{sec:ch16-shadow}
\index{shadowing}

\noindent
So every long chaotic computation in this book followed the wrong orbit.
Did it draw the wrong picture? There is a way to find out, and it rests on
one observation.

\begin{think}
Running the logistic map forwards doubles a small error at every step, on
average. What do you expect running it backwards, from each value to the one
before it, to do to a small error?
\end{think}
\reveal{Halve it, on average. Whatever a forward step stretches, the backward
step squeezes.}
So start at the end of a computed orbit and walk back. Two numbers lead to
any value in one step, one below $\tfrac{1}{2}$ and one above it; at each
step take the one on the same side as the computed value. Because walking
back shrinks errors, the walk stays close to the computed orbit all the way
to the start, and what it finds, worked to four hundred digits, is a genuine
orbit of the rule with no rounding in it.

\pyregion{code/ch16/shadow.py}{back}{Walking a computed orbit backwards to find the true orbit it follows}{lst:ch16-back}

\transcript{ch16-shadow}

\noindent
Read the first three lines. The computed orbit of $\num{0.1}$ left the truth
after \val{ch16.wrong.a} steps. But there is a start, differing from one
tenth by about \val{ch16.shadow.start}, whose exact orbit stays within
\val{ch16.shadow.gap} of every one of the \val{ch16.shadow.steps} computed
values. The computer did not follow the orbit it was asked for. It followed,
faithfully, the orbit of a start so close to it that no measurement could
tell the two apart. That true orbit is called a \emph{shadow}.

\begin{rigourbox}
The \emph{shadowing lemma}, proved by Dmitri Anosov and Rufus Bowen for the
tidiest kind of chaotic system, says this. Suppose a system stretches in some
directions and squeezes in others at every point, at rates that never weaken
(mathematicians call it \emph{uniformly hyperbolic}). Then however closely you
want a computed orbit to be shadowed, there is a size of error per step small
enough to guarantee it: any sequence that makes errors no bigger than that at
each step stays that close to some true orbit, for ever. The doubling map is
of this kind. The logistic map at $r = 4$ is not quite: near
$x = \tfrac{1}{2}$ its curve is flat, and there the backward walk stretches
instead of squeezing. For it, shadowing has been confirmed by careful
computation over very long but finite stretches, notably by Hammel, Yorke and
Grebogi in 1987, rather than proved for ever. This book proves neither; the
proofs live in graduate texts on dynamical systems.
\end{rigourbox}

\noindent
Now the payoff. A shadow is a true orbit, so it spends its time the way true
orbits of this map do, and the last two lines of the output measure it: over
a million steps, both formulas spend the fraction \val{ch16.stats.a} of their
time below $\num{0.1}$, and the exact answer, known from a formula this book
does not derive, is \val{ch16.stats.exact}. Two programs whose paths parted
after fifty steps agree about where the orbit lives. That is why the lobe
fractions of Chapter~\ref{ch:lorenz} and the climate statistics of
Chapter~\ref{ch:weather} can be trusted when no single computed path can.

\plotfig{ch16-statistics}{Left: where two double-precision runs of a million
steps spend their time, against the exact answer. Right: the share of steps
below \num{0.1} as the runs lengthen. Both double runs settle on the exact
value; the single-precision run, caught in a loop, settles on a wrong
one.}{statistics that survive a wrong path}

\mermaidfig{ch16-shadow}{A computed orbit is not the orbit of its own start,
but it stays close to the true orbit of a nearby start, and that orbit's
statistics are the ones the picture shows.}{a computed orbit and its shadow}

\section{Machines that must repeat}
\label{sec:ch16-random}
\index{random numbers}\index{linear congruential generator}

\noindent
There is one more way a computer differs from mathematics, and it is the
plainest. Between $0$ and $1$ there are \val{ch16.count.single}
single-precision numbers and \val{ch16.count.double} doubles. Enormous, and
finite. So an orbit on a computer must sooner or later land on a value it has
met before, and since the rule always gives the same next value from the same
value, it then goes round the same loop for ever.

In single precision the loops are short enough to find. From $\num{0.1}$ the
logistic map at $r = 4$ wanders for \val{ch16.loop.tail} steps and then goes
round a loop of \val{ch16.loop.length} values for ever; from
\val{ch16.loop.starts} different starts it ends in only \val{ch16.loop.kinds}
different loops. The true map from almost any start never repeats at all. The
loop also spoils the statistics that saved the day in the last section: the
single-precision run spends the fraction \val{ch16.stats.single} of its time
below $\num{0.1}$, not \val{ch16.stats.exact}, as the right half of the last
plot shows. Doubles have loops too, but they are far too long to reach, and
that is the practical reason to compute chaos in double precision.

\begin{think}
The logistic map at $r = 4$ never repeats, spreads its values over the whole
of $0$ to $1$, and looks random at a glance. Would it make a good
random-number generator?
\end{think}
\reveal{No, for three reasons this chapter has measured. Its values pile up
near $0$ and $1$ \dash{} about a fifth of its time is spent below
$\num{0.1}$, not a tenth. Each value is fixed by the one before it. And on a
computer it falls into a loop.}

\begin{trapbox}
\textbf{\enquote{A chaotic map makes a good random-number generator.}} Chaos
makes the future unpredictable to someone who does not know the start
exactly, which is part of what randomness needs and not all of it. A good
generator must also spread its values evenly, show no pattern between one
value and the next, and repeat only after a loop so long that nobody could
ever go round it. Those properties are designed with number theory and
checked with batteries of statistical tests; they do not come free with
chaos.
\end{trapbox}

\noindent
A random-number generator in a computer is a deterministic rule on a finite
set of values, built on purpose to have one enormous loop. The oldest kind,
the \emph{linear congruential generator}, is one line: multiply by a
number~$a$, add a number~$c$, and keep the remainder after dividing by a
number~$m$. With $a = 5$, $c = 3$ and $m = 16$, starting from $0$, you get
$3$, then $5 \times 3 + 3 = 18$, remainder $2$, then $13$, $4$, $7$, $6$, $1$,
$8$, $11$, $10$, $5$, $12$, $15$, $14$, $9$ and back to $0$: all sixteen
possible values before a single repeat, which is the best a machine with
sixteen values can do. The doubling map is a generator of this kind with
$a = 2$ and $c = 0$, and on a machine with $k$ binary digits it dies in $k$
steps, which is the collapse of the second section in a new costume.

Real generators use the same idea on a grand scale. Python's own, the
Mersenne Twister of Makoto Matsumoto and Takuji Nishimura (1998), goes round a
loop of $2^{19937} - 1$ values. Because it is deterministic, giving it the
same starting value, called the \emph{seed}, gives the same
\enquote{random} numbers every time, which is how a scientist makes a random
experiment repeatable.

\begin{lab}{l16_03_loops}{Every orbit ends in a loop}
Write a linear congruential generator and a function that finds, for any rule
and start, how many steps the orbit takes to enter its loop and how long the
loop is. Check it on the sixteen-value generator above and on the doubling
map with eight binary digits, then use it to find the loop of the
single-precision logistic map for yourself.
\end{lab}

\noindent
The practical lesson of the chapter fits in a sentence: never trust a single
long run of a chaotic computation. Run it again with the formula rearranged,
with more precision or less, with a different step size, and believe what
survives. The path past fifty steps will not survive, and was never the
answer. The statistics will, and are.

\begin{summarybox}
\begin{itemize}
\item A computer stores a number as a \textbf{double}: \val{ch16.bits.double}
  significant binary digits, about \val{ch16.decimal.double} decimal ones.
  Most decimals, $\num{0.1}$ among them, cannot be stored exactly, and every
  sum and product is \textbf{rounded} again.
\item The \textbf{doubling map} reads out its start one binary digit at a
  time, so on a computer every orbit dies at $0$ once the digits run out
  \dash{} from $\num{0.1}$ at step \val{ch16.collapse.tenth}, while the true
  orbit loops for ever. The arithmetic is exact; the storable starts are the
  rare ones that die.
\item Two programs for the same rule round differently and part company. At
  $r = 4$ the logistic map uses up \textbf{one binary digit per step}, so a
  double lasts about fifty steps and a single about twenty.
\item \textbf{Exact arithmetic} doubles its digits every step and cannot
  finish; \textbf{more digits} buy a fixed number of steps each, never an
  exact orbit.
\item A computed orbit is \textbf{shadowed} by a true orbit from a nearby
  start, so its \textbf{statistics} are right even when its path is wrong.
\item A computer has finitely many numbers, so every computed orbit ends in a
  \textbf{loop}. A random-number generator is a deterministic rule built to
  make that loop enormous; a chaotic map is not a good one.
\end{itemize}
\end{summarybox}

\canyou

\begin{checkpoint}
\item Why does $\num{0.1} + \num{0.2}$ not equal $\num{0.3}$ on a computer?
  \answerto{Neither $\num{0.1}$ nor $\num{0.2}$ has an exact binary form, so
  each is stored as the nearest double, and the sum is rounded once more. The
  result is a double just above the one stored for $\num{0.3}$.}
\item The computer stores $\num{0.3}$ as a whole number divided by
  $2^{\val{ch16.collapse.three}}$. After how many steps of the doubling map is
  its orbit exactly $0$, and what does the true orbit of $\tfrac{3}{10}$ do?
  \answerto{After \val{ch16.collapse.three} steps, one for each binary digit.
  The true orbit goes $\tfrac{3}{5}$, $\tfrac{1}{5}$, $\tfrac{2}{5}$,
  $\tfrac{4}{5}$, $\tfrac{3}{5}$ and round that loop of four for ever.}
\item A program runs the logistic map at $r = 4$ in single precision, which
  keeps \val{ch16.bits.single} binary digits. Roughly how many steps is it
  good for, and why?
  \answerto{About twenty (measured: \val{ch16.wrong.single}). Its rounding
  error is about one part in $2^{24}$, and the map doubles an error every
  step, using up one binary digit per step.}
\item Run exactly from $\tfrac{1}{10}$, the bottom of the fraction has
  \val{ch16.exact.ten} digits after ten steps. About how many does it have
  after twenty?
  \answerto{About \val{ch16.exact.twenty}: the number of digits doubles every
  step, and ten more doublings multiply it by about a thousand.}
\item After a million steps of a chaotic simulation, which can you trust: the
  value at the last step, or the fraction of time spent below $\num{0.1}$?
  Why?
  \answerto{The fraction of time. The path has followed a different orbit
  since about step fifty, but it stays close to a true orbit from a nearby
  start (a shadow), and that orbit spends its time as true orbits do.}
\item Why must every orbit computed on a machine eventually repeat, and what
  does a good random-number generator do about it?
  \answerto{The machine has only finitely many numbers, so the orbit must
  revisit a value, and the rule then repeats everything after it. A good
  generator is designed so that its loop is too long ever to be gone round.}
\end{checkpoint}
